\documentclass[11pt]{article}
\usepackage{amsmath,amssymb,amsthm,amsfonts}
\usepackage{geometry}
\geometry{margin=1in}
\usepackage{algorithm}
\usepackage{algpseudocode}

\newtheorem{theorem}{Theorem}
\newtheorem{lemma}{Lemma}
\newtheorem{assumption}{Assumption}
\newtheorem{definition}{Definition}
\newtheorem{corollary}{Corollary}

\title{Stochastic Subgradient (Stochastic Gradient) Descent:\\
Algorithm and Convergence Analysis under Non-convex Smooth Objectives}
\author{}
\date{}

\begin{document}
\maketitle

%=====================================================================
\section*{Algorithm Name}
\textbf{Stochastic (Sub)Gradient Descent (SGD).}
The algorithm iteratively updates parameters by sampling a single random
data point, computing a stochastic (sub)gradient, and stepping in the
negative gradient direction. Although the algorithm is classically stated
for convex Lipschitz objectives, here we analyze it under the stated
assumptions of \emph{non-convex, smooth} objective functions, where the
appropriate convergence criterion is the expected squared gradient norm.

%=====================================================================
\section*{Mathematical Formulation}
Let $f:\mathbb{R}^d\to\mathbb{R}$ be the objective. We assume access to a
stochastic oracle returning $g_t=\nabla f(w_t;z_t)$ where $z_t$ is a random
sample drawn i.i.d.\ at iteration $t$.

\textbf{Update rule:}
\[
w_{t+1} = w_t - \eta_t\, g_t,\qquad g_t=\nabla f(w_t;z_t).
\]

\textbf{Hyperparameters:}
\begin{itemize}
  \item Step size (learning rate) $\eta_t>0$; either constant ($\eta_t=\eta$)
        or diminishing (e.g.\ $\eta_t=c/\sqrt{t}$).
  \item Total number of iterations $T$.
  \item Initialization $w_0\in\mathbb{R}^d$.
\end{itemize}

\textbf{Oracle conditions:}
\[
\mathbb{E}[g_t\mid w_t]=\nabla f(w_t),\qquad
\mathbb{E}\!\left[\|g_t-\nabla f(w_t)\|^2\mid w_t\right]\le\sigma^2.
\]

%=====================================================================
\section*{Pseudocode}
\begin{algorithm}[H]
\caption{Stochastic Gradient Descent (SGD)}
\begin{algorithmic}[1]
\State \textbf{Input:} initial point $w_0$, step sizes $\{\eta_t\}_{t=0}^{T-1}$, iterations $T$
\For{$t=0,1,2,\dots,T-1$}
    \State Sample a data point $z_t$ i.i.d.\ from the data distribution
    \State Compute stochastic gradient $g_t=\nabla f(w_t;z_t)$
    \State Update $w_{t+1}\gets w_t-\eta_t\,g_t$
\EndFor
\State \textbf{Termination:} stop when $t=T$ or $\|g_t\|\le\texttt{tol}$
\State \textbf{Output:} $w_R$ where $R$ is sampled with probability proportional to $\eta_R$ (or $\bar w_T=\frac{1}{T}\sum_t w_t$ in the convex case)
\end{algorithmic}
\end{algorithm}

%=====================================================================
\section*{Dry Run Example}
Consider $f(w)=(w-3)^2$, with $w_0=0$, $\eta=0.1$, deterministic gradient
$\nabla f(w)=2(w-3)$. We perform $3$ iterations.

\textbf{Iteration 1 ($t=0$):}
\begin{align*}
g_0 &= 2(w_0-3)=2(0-3)=-6,\\
w_1 &= w_0-\eta g_0 = 0-0.1(-6)=0.6,\\
f(w_1) &= (0.6-3)^2=(-2.4)^2=5.76.
\end{align*}

\textbf{Iteration 2 ($t=1$):}
\begin{align*}
g_1 &= 2(w_1-3)=2(0.6-3)=-4.8,\\
w_2 &= w_1-\eta g_1 = 0.6-0.1(-4.8)=1.08,\\
f(w_2) &= (1.08-3)^2=(-1.92)^2=3.6864.
\end{align*}

\textbf{Iteration 3 ($t=2$):}
\begin{align*}
g_2 &= 2(w_2-3)=2(1.08-3)=-3.84,\\
w_3 &= w_2-\eta g_2 = 1.08-0.1(-3.84)=1.464,\\
f(w_3) &= (1.464-3)^2=(-1.536)^2=2.359296.
\end{align*}

The objective decreases monotonically: $9 \to 5.76 \to 3.6864 \to 2.3593$,
converging toward the minimizer $w^*=3$, $f(w^*)=0$.

%=====================================================================
\section*{Assumptions}
\begin{assumption}[Lower boundedness]\label{as:bounded}
$f$ is bounded below: $f^\star:=\inf_{w}f(w)>-\infty$.
\end{assumption}

\begin{assumption}[$L$-smoothness]\label{as:smooth}
$f$ is differentiable and its gradient is $L$-Lipschitz:
\[
\|\nabla f(w)-\nabla f(v)\|\le L\|w-v\|,\qquad\forall w,v\in\mathbb{R}^d.
\]
Equivalently (descent lemma),
\[
f(v)\le f(w)+\langle\nabla f(w),v-w\rangle+\tfrac{L}{2}\|v-w\|^2.
\]
\end{assumption}

\begin{assumption}[Unbiased stochastic gradients]\label{as:unbiased}
$\mathbb{E}[g_t\mid w_t]=\nabla f(w_t)$.
\end{assumption}

\begin{assumption}[Bounded variance]\label{as:var}
$\mathbb{E}\!\left[\|g_t-\nabla f(w_t)\|^2\mid w_t\right]\le\sigma^2$.
\end{assumption}

Note: under non-convexity the natural optimality measure is the squared
gradient norm $\|\nabla f(w)\|^2$ (stationarity), not the function-value gap.

%=====================================================================
\section*{Main Convergence Theorem}
\begin{theorem}[Constant step size]\label{thm:const}
Under Assumptions~\ref{as:bounded}--\ref{as:var}, run SGD with constant
step size $\eta_t=\eta\le 1/L$. Then
\[
\frac{1}{T}\sum_{t=0}^{T-1}\mathbb{E}\!\left[\|\nabla f(w_t)\|^2\right]
\le \frac{2\bigl(f(w_0)-f^\star\bigr)}{\eta T}+L\eta\sigma^2.
\]
\end{theorem}

\begin{corollary}[Optimized step size]\label{cor:opt}
Choosing
$\eta=\min\!\left\{\dfrac{1}{L},\,\sqrt{\dfrac{2(f(w_0)-f^\star)}{L\sigma^2 T}}\right\}$
yields
\[
\frac{1}{T}\sum_{t=0}^{T-1}\mathbb{E}\!\left[\|\nabla f(w_t)\|^2\right]
= O\!\left(\frac{1}{\sqrt{T}}\right).
\]
\end{corollary}

%=====================================================================
\section*{Theoretical Properties}
\begin{itemize}
  \item \textbf{Stationarity rate:} $O(1/\sqrt{T})$ for the average squared
        gradient norm, matching the known lower bound for first-order
        stochastic methods on smooth non-convex problems.
  \item \textbf{Stability:} requires $\eta\le1/L$; larger steps may diverge
        because the descent lemma no longer guarantees decrease.
  \item \textbf{Hyperparameter sensitivity:} a bias term $L\eta\sigma^2$
        persists for constant $\eta$; diminishing $\eta_t$ removes the bias
        at the cost of a slower transient.
  \item \textbf{Baseline comparison:} in the deterministic case ($\sigma=0$),
        the bound reduces to $O(1/T)$, recovering classical gradient descent.
\end{itemize}

%=====================================================================
\section*{Discussion}
Although SGD is classically motivated for convex Lipschitz objectives, the
same single-sample subgradient update admits a clean analysis under smooth
non-convex objectives. The key change is the optimality criterion: instead
of bounding $f(\bar w_T)-f^\star$, we bound the average squared gradient
norm, the standard stationarity measure. The variance term $L\eta\sigma^2$
captures the price of stochasticity.

%=====================================================================
\section*{Preliminaries (Definitions and Lemmas)}
\begin{lemma}[Descent Lemma]\label{lem:descent}
If $f$ is $L$-smooth, then for all $w,v$:
\[
f(v)\le f(w)+\langle\nabla f(w),v-w\rangle+\frac{L}{2}\|v-w\|^2.
\]
\end{lemma}
\begin{proof}
By the fundamental theorem of calculus,
$f(v)-f(w)=\int_0^1\langle\nabla f(w+s(v-w)),v-w\rangle\,ds$.
Subtract $\langle\nabla f(w),v-w\rangle$:
\[
f(v)-f(w)-\langle\nabla f(w),v-w\rangle
=\int_0^1\langle\nabla f(w+s(v-w))-\nabla f(w),\,v-w\rangle\,ds.
\]
By Cauchy--Schwarz and $L$-Lipschitzness,
$\langle\nabla f(w+s(v-w))-\nabla f(w),v-w\rangle\le Ls\|v-w\|^2$.
Integrating $\int_0^1 Ls\,ds=L/2$ gives the result.
\end{proof}

%=====================================================================
\section*{Main Proof (Step-by-Step Derivation)}
\begin{proof}[Proof of Theorem~\ref{thm:const}]
\textbf{[Step 1]} Apply the Descent Lemma~\ref{lem:descent} with
$w=w_t$ and $v=w_{t+1}=w_t-\eta g_t$:
\[
f(w_{t+1})\le f(w_t)+\langle\nabla f(w_t),\,w_{t+1}-w_t\rangle
+\frac{L}{2}\|w_{t+1}-w_t\|^2.
\]

\textbf{[Step 2]} Substitute $w_{t+1}-w_t=-\eta g_t$:
\[
f(w_{t+1})\le f(w_t)-\eta\langle\nabla f(w_t),g_t\rangle
+\frac{L\eta^2}{2}\|g_t\|^2.
\]

\textbf{[Step 3]} Take conditional expectation given $w_t$. Using
Assumption~\ref{as:unbiased}, $\mathbb{E}[g_t\mid w_t]=\nabla f(w_t)$, so
\[
\mathbb{E}[\langle\nabla f(w_t),g_t\rangle\mid w_t]
=\langle\nabla f(w_t),\nabla f(w_t)\rangle=\|\nabla f(w_t)\|^2.
\]
Hence
\[
\mathbb{E}[f(w_{t+1})\mid w_t]\le f(w_t)-\eta\|\nabla f(w_t)\|^2
+\frac{L\eta^2}{2}\,\mathbb{E}[\|g_t\|^2\mid w_t].
\]

\textbf{[Step 4]} Decompose the second moment. Since the cross term
vanishes by unbiasedness:
\begin{align*}
\mathbb{E}[\|g_t\|^2\mid w_t]
&=\mathbb{E}[\|g_t-\nabla f(w_t)+\nabla f(w_t)\|^2\mid w_t]\\
&=\mathbb{E}[\|g_t-\nabla f(w_t)\|^2\mid w_t]
+\|\nabla f(w_t)\|^2\\
&\quad+2\,\mathbb{E}[\langle g_t-\nabla f(w_t),\nabla f(w_t)\rangle\mid w_t].
\end{align*}
The cross term is $2\langle\mathbb{E}[g_t\mid w_t]-\nabla f(w_t),\nabla f(w_t)\rangle=0$.
By Assumption~\ref{as:var},
\[
\mathbb{E}[\|g_t\|^2\mid w_t]\le\|\nabla f(w_t)\|^2+\sigma^2.
\]

\textbf{[Step 5]} Substitute into the Step 3 bound:
\[
\mathbb{E}[f(w_{t+1})\mid w_t]
\le f(w_t)-\eta\|\nabla f(w_t)\|^2
+\frac{L\eta^2}{2}\bigl(\|\nabla f(w_t)\|^2+\sigma^2\bigr).
\]

\textbf{[Step 6]} Group the gradient-norm terms:
\[
\mathbb{E}[f(w_{t+1})\mid w_t]
\le f(w_t)-\eta\left(1-\frac{L\eta}{2}\right)\|\nabla f(w_t)\|^2
+\frac{L\eta^2\sigma^2}{2}.
\]

\textbf{[Step 7]} Use the step-size condition $\eta\le 1/L$, which gives
$\frac{L\eta}{2}\le\frac12$, hence $1-\frac{L\eta}{2}\ge\frac12$. Therefore
\[
\mathbb{E}[f(w_{t+1})\mid w_t]
\le f(w_t)-\frac{\eta}{2}\|\nabla f(w_t)\|^2
+\frac{L\eta^2\sigma^2}{2}.
\]

\textbf{[Step 8]} Take total expectation (tower property) over all
randomness up to $t$:
\[
\mathbb{E}[f(w_{t+1})]
\le \mathbb{E}[f(w_t)]-\frac{\eta}{2}\,\mathbb{E}[\|\nabla f(w_t)\|^2]
+\frac{L\eta^2\sigma^2}{2}.
\]

\textbf{[Step 9]} Rearrange to isolate the gradient term:
\[
\frac{\eta}{2}\,\mathbb{E}[\|\nabla f(w_t)\|^2]
\le \mathbb{E}[f(w_t)]-\mathbb{E}[f(w_{t+1})]
+\frac{L\eta^2\sigma^2}{2}.
\]

\textbf{[Step 10]} Sum over $t=0,1,\dots,T-1$ (telescoping):
\[
\frac{\eta}{2}\sum_{t=0}^{T-1}\mathbb{E}[\|\nabla f(w_t)\|^2]
\le \mathbb{E}[f(w_0)]-\mathbb{E}[f(w_T)]
+\frac{L\eta^2\sigma^2}{2}T.
\]

\textbf{[Step 11]} Since $f(w_T)\ge f^\star$ (Assumption~\ref{as:bounded})
and $f(w_0)$ is deterministic,
\[
\frac{\eta}{2}\sum_{t=0}^{T-1}\mathbb{E}[\|\nabla f(w_t)\|^2]
\le f(w_0)-f^\star+\frac{L\eta^2\sigma^2}{2}T.
\]

\textbf{[Step 12]} Multiply both sides by $\frac{2}{\eta T}$:
\[
\frac{1}{T}\sum_{t=0}^{T-1}\mathbb{E}[\|\nabla f(w_t)\|^2]
\le \frac{2\bigl(f(w_0)-f^\star\bigr)}{\eta T}+L\eta\sigma^2.
\]
This is exactly the claimed bound in Theorem~\ref{thm:const}.
\end{proof}

%=====================================================================
\section*{Rate Derivation (With Constants)}
\begin{proof}[Proof of Corollary~\ref{cor:opt}]
\textbf{[Step 13]} Denote $\Delta_0:=f(w_0)-f^\star$. The bound from
Theorem~\ref{thm:const} is
\[
\Phi(\eta):=\frac{2\Delta_0}{\eta T}+L\eta\sigma^2.
\]

\textbf{[Step 14]} Ignoring the constraint $\eta\le1/L$, minimize over
$\eta>0$. Setting $\Phi'(\eta)=-\frac{2\Delta_0}{\eta^2 T}+L\sigma^2=0$:
\[
\eta^\star=\sqrt{\frac{2\Delta_0}{L\sigma^2 T}}.
\]

\textbf{[Step 15]} Substitute $\eta^\star$ back:
\[
\Phi(\eta^\star)
=\frac{2\Delta_0}{T}\sqrt{\frac{L\sigma^2 T}{2\Delta_0}}
+L\sigma^2\sqrt{\frac{2\Delta_0}{L\sigma^2 T}}
=2\sqrt{\frac{2\Delta_0 L\sigma^2}{T}}.
\]
Both terms equal $\sqrt{2\Delta_0 L\sigma^2/T}$, summing to the above.

\textbf{[Step 16]} To respect $\eta\le1/L$ we take
$\eta=\min\{1/L,\eta^\star\}$. For large $T$, $\eta^\star\le1/L$, so
\[
\frac{1}{T}\sum_{t=0}^{T-1}\mathbb{E}[\|\nabla f(w_t)\|^2]
\le 2\sqrt{\frac{2L\sigma^2\Delta_0}{T}}=O\!\left(\frac{1}{\sqrt{T}}\right).
\]
For small $T$ (when $\eta=1/L$ binds), the bound becomes
$\frac{2L\Delta_0}{T}+\sigma^2$, dominated by the $O(1/\sqrt{T})$ regime
asymptotically.
\end{proof}

%=====================================================================
\section*{Special Cases}
\begin{corollary}[Deterministic gradient descent, $\sigma=0$]
If gradients are exact ($\sigma=0$) and $\eta=1/L$, then
\[
\frac{1}{T}\sum_{t=0}^{T-1}\|\nabla f(w_t)\|^2
\le\frac{2L\bigl(f(w_0)-f^\star\bigr)}{T}=O\!\left(\frac{1}{T}\right),
\]
recovering the classical non-convex GD stationarity rate.
\end{corollary}

\begin{corollary}[Diminishing step size]
If $\eta_t=c/\sqrt{t+1}$ with $c\le1/L$, repeating Steps 8--11 with
non-constant $\eta_t$ and using $\sum_{t=0}^{T-1}\eta_t=\Theta(\sqrt{T})$
and $\sum_{t=0}^{T-1}\eta_t^2=\Theta(\log T)$ yields
\[
\min_{0\le t<T}\mathbb{E}[\|\nabla f(w_t)\|^2]
\le \frac{\,f(w_0)-f^\star+\tfrac{L\sigma^2 c^2}{2}\log T\,}
{c\,(\sqrt{T}-1)}
=O\!\left(\frac{\log T}{\sqrt{T}}\right),
\]
removing the persistent variance bias at the cost of a logarithmic factor.
\end{corollary}

\end{document}